\section{Introducción}

La metrología es la ciencia que establece las bases de la medición y garantiza que
los resultados obtenidos en el laboratorio o en la industria sean comparables,
confiables y trazables a patrones de referencia. En este contexto, la calibración
de instrumentos de medición constituye el procedimiento mediante el cual se
comprueba su comportamiento frente a un patrón de mayor exactitud, se cuantifican
sus errores de indicación y se estima la incertidumbre asociada conforme a la
guía GUM~\parencite{jcgm2008gum}.

El presente informe documenta las actividades de calibración desarrolladas durante
las prácticas pre-profesionales realizadas en Physics Technology Investigation
(PHYTECHIN EIRL). El trabajo se centró en tres instrumentos: un pie de rey
Mitutoyo 530-115 y un comparador digital de cuadrante Mitutoyo ID-S1012M, ambos
calibrados por comparación directa con bloques patrón siguiendo los procedimientos
normalizados del INACAL~\parencite{inacal_pc012,inacal_pc014} y la norma
ISO~3650~\parencite{iso3650_1998}, y un medidor de espesores por ultrasonido
INSIZE ISU-100D calibrado con un patrón escalonado de acero según el procedimiento
PC-029~\parencite{inacal_pc029}.

El análisis se apoya en los fundamentos de probabilidad y en la ley de propagación
de incertidumbres que se desarrollan en el marco teórico. Las componentes de
incertidumbre tipo~A se obtuvieron del tratamiento estadístico de lecturas
repetidas, en tanto que las componentes tipo~B se evaluaron a partir de
certificados de calibración, especificaciones del fabricante y límites conocidos,
empleando la distribución de probabilidad adecuada a cada caso. Como resultado se
obtuvieron las expresiones de la incertidumbre combinada y expandida, así como la
declaración de conformidad del pie de rey según la regla de decisión de la guía
ILAC-G8~\parencite{ilac2019decision}.
